\section{Introduction}\label{sec:intro}

Social accounting matrices organise an economy's transactions---production,
factor incomes, institutional transfers, accumulation, and external
flows---into a single square array in which every account's receipts equal
its payments. They are the standard database for multiplier analysis and
computable general equilibrium modelling, and in many countries they are the
principal quantitative bridge between official statistics and policy
simulation.

Yet SAM construction has a reproducibility problem. The way matrices are
built has changed remarkably little since \citet{pyattround1985}:
heterogeneous official tables are reconciled through long chains of
spreadsheet operations, expert judgments, and balancing passes that are
rarely documented at the level where they act---the individual cell. The
consequences are familiar to every practitioner: published SAMs are trusted
and used, but they can seldom be reproduced, audited, or updated by anyone
other than their authors; knowledge of how a matrix was actually built
leaves when its builder does; and each new country or benchmark year
restarts the work nearly from scratch. Section~\ref{sec:related} argues
that this is a property of the prevailing \emph{medium}---prose describing
spreadsheets---rather than of any particular team.

This paper demonstrates that the problem is unnecessary: SAM construction
can move from an opaque artisanal exercise to an auditable, executable
statistical-production process. We introduce a construction and audit
framework with four commitments: every source file is fingerprinted and
registered; every classification, concordance, and cell formula is
machine-readable and reviewable; every balancing adjustment is logged at
the cell where it acts; and the accounting identities that define a SAM
are enforced as executable checks rather than assumed. We then apply the
framework end to end, building a complete 2019 SAM for South Africa from
publicly obtainable official sources and access-controlled survey
microdata---inputs available to any researcher at no cost, with no
privileged access\footnote{``Publicly obtainable'' means: published
statistical releases plus survey microdata obtainable by any registered
researcher under DataFirst's standard access conditions. The one input
that is not obtainable at all---the benchmark's unpublished customs
figure---is treated separately and priced
(Section~\ref{sec:results:macro}); the \emph{public-proxy} scenario of
our SAM uses published sources only.}---and auditing it, cell by cell,
against the strongest available benchmark: the 2019 South African SAM
published by UNU-WIDER \citep{vanseventer2023}, which is unusually well
documented by the standards of the genre and ships with a technical note
mapping its macro structure to South African Reserve Bank (SARB)
national-accounts series.

South Africa is the evidence, not the point. The audit produces three kinds
of result that we argue generalise. First, \emph{exact replication where
the sources allow it}: the benchmark's entire production and commodity
side reproduces from the official supply and use tables to within
R0.5~million on all 61 activities and all commodities, and 43 of 45
computable institutional macro cells match the SARB series exactly
(Table~\ref{tab:headline} summarises). Second, \emph{numerical recovery of
undocumented choices}: where the benchmark's documentation is silent or
mistaken, the data identify what was actually done---the edition and
measure behind its capital-stock split, the unique sign structure of a
misprinted aggregation formula, and the household grouping variable.
Third, \emph{a priced replicability boundary}: the benchmark rests on
exactly one non-public input, and substituting its public proxy changes
import duties by $+16.1$ per cent and product taxes by $-1.5$ per cent; a
small unrecorded balancing adjustment and a handful of sector-level share
edits are likewise isolated and measured---each detected because it left
an arithmetic trace, which is the limit of what this method can claim.

Two further applications then place South Africa on a spectrum, and the
three-country design is deliberately asymmetric: one deep audit, one
data-boundary test, one transfer-and-generation demonstration. The
applications are not three comparable country studies---each answers a
different question about the same framework under a different publication
regime, and their unequal depth is the research design, not uneven
execution. Applied to Kenya---where the machine-readable supply--use
detail behind IFPRI's widely used 2019 SAM is not published---the
framework verifies the matrix's structure and macro controls against the
national accounts and locates exactly where verification must stop: the
published accounts carry a statistical discrepancy of 3.6~per~cent of
GDP, a balanced SAM cannot carry one, and the deviation pattern is
consistent with the discrepancy being absorbed primarily through the
consumption and trade aggregates. Applied to the Netherlands---where
harmonised tables are a public API---the framework needs no benchmark at
all: three API calls, a 70-line reader, and one script generate a 2019
macro accounting matrix whose inputs pass every identity exactly and
whose residuals are each attributed to a named ESA boundary item, and a
logged balancing step then closes it into a macro SAM
(Section~\ref{sec:applications}).

Our contribution is therefore not a new balancing estimator. That
literature is active rather than closed---cross-entropy estimation
\citep{robinson2001}, GRAS for matrices with negative entries
\citep{junius2003,temurshoev2013}, and KRAS for conflicting and
unequally reliable constraints \citep{lenzen2009} each address a
different compilation problem---and this paper adds nothing to it. What
we contribute instead is a demonstration, on a full-scale
national SAM, of five things at once: a reproducible construction pipeline;
an audit methodology whose discrepancies are detected by
\emph{deterministic reconciliation against accounting identities} rather
than by reading documentation; a reconstruction of a published national SAM
from publicly obtainable sources with every departure disclosed and
quantified; a
taxonomy of the ways documentation and data diverge in even the
best-documented matrices; and an open toolkit---readers for supply--use,
SAM, central-bank, and survey sources; concordance and account-taxonomy
machinery; RAS balancing with mandatory diagnostics; a generator for
workbooks whose totals and balance checks are live formulas---offered as
reproducible infrastructure for official SAM construction and released
with the replication package.\footnote{The toolkit (\texttt{edikit}) and
all scripts, registries, and derived tables are part of the replication
package; see the Reproducibility and AI disclosure section.} The framework
doubles as a practical construction workflow: the pipeline generates and
balances the matrix mechanically, so that scarce expert time is spent only
on the residual cells that identities cannot determine
(Section~\ref{sec:discussion:workflow}).

The idea the paper keeps returning to is simple: \emph{cell-level
provenance and machine-readable construction rules, published alongside
every SAM, would make such audits---and such SAMs---routine.} The paper
proceeds as follows. Section~\ref{sec:related} situates the work and
contrasts the prevailing construction workflow with the framework's.
Section~\ref{sec:data} describes the source registry.
Section~\ref{sec:method} sets out the pipeline and the audit method.
Section~\ref{sec:results} reports the South African replication block by
block. Section~\ref{sec:applications} presents the Kenyan and Dutch
applications and the replicability spectrum they trace.
Section~\ref{sec:discussion} draws out the taxonomy, the construction
workflow, and implications for statistical practice, and
Section~\ref{sec:conclusion} concludes.
